\subsection{Eligibility, screening and study quality}
Include general synthetic-data surveys and foundational generation methods
as background for RQ1. For the focused stage, include studies generating,
simulating or translating automotive LiDAR point clouds, and studies
evaluating their use for detection or semantic segmentation. Tag
completion and novel-view reconstruction separately from new-scene generation.
Do not exclude a generation paper merely because it lacks downstream
experiments; record that absence as a limit on its utility evidence.

Exclude unrelated applications and object-only shape synthesis from the
focused synthesis unless their connection to driving scans is demonstrated.
Exclude duplicate versions and sources without enough methodological detail
to interpret their contribution. Link preprints to published versions and
avoid counting the same study twice.

Deduplicate by DOI and title, then screen titles and abstracts before
full texts. Each author records an independent decision and reason;
discuss disagreements and preserve the agreed decision alongside them.
Report actual counts and full-text exclusion reasons.
Assess generation-task relevance, sensor and annotation documentation,
independent real test data, comparable baselines, uncertainty and
reproducibility. Use these judgements to qualify the synthesis.
